\textbf{Background/Objectives:} Small white matter lesions are easily overlooked on MRI in multiple sclerosis (MS), yet the 2024 McDonald criteria and MAGNIMS--CMSC--NAIMS recommendations consider lesions from 3~mm in size, and new lesions on follow-up MRI inform treatment decisions. Segmentation models are usually trained to maximize voxel overlap (Dice score), which can remain high while small lesions are missed. We evaluated whether a lesion-level training objective improves the detection of small MS lesions. \textbf{Methods:} CATMIL adds two lesion-level terms to the standard nnU-Net training loss: one gives every lesion equal weight regardless of its volume, and one penalizes any lesion that is left completely undetected. Models using T1-weighted, T2-weighted, and FLAIR images were trained on two public datasets (MSLesSeg, 53 patients; 3D-MR-MS, 30 patients), evaluated on held-out patients, and compared with four alternative training losses using patient-level bootstrap confidence intervals. \textbf{Results:} For lesions of at least 3~mm, CATMIL detected 94.4\% of reference lesions versus 87.0\% with the standard loss on MSLesSeg (difference +7.4 percentage points, 95\% CI +3.2 to +12.9) and 89.4\% versus 83.0\% on 3D-MR-MS. Completely missed lesions decreased by about 48\% (2.6 vs.\ 5.0 per scan) with similar Dice overlap, and the reduction was reproduced in the 47 cross-validation patients (4.1 vs.\ 7.5 per scan). The cost was additional small false-positive detections, most of which were removed by a simple size filter. \textbf{Conclusions:} Lesion-level training substantially reduces missed small MS lesions at the clinically relevant 3~mm size, making it a candidate second reader that flags lesions for radiologist confirmation.
